\ifx\allfiles\undefined
\documentlecture[12pt, a4paper, oneside, UTF8]{ctexbook}  %  这一句是新增加的
\usepackage[dvipsnames]{xcolor}
\usepackage{amsmath}   % 数学公式
\usepackage{graphicx}
\usetikzlibrary{arrows, calc, decorations.pathmorphing}
\allowdisplaybreaks % 允许公式跨页换行
\newcommand{\pa}{\partial}
\newcommand{\mT}{\raisebox{0.1ex}{$\scriptstyle -T$}} % 调整高度为 0.1ex
\newcommand{\lmT}{\raisebox{-0.85ex}{$\scriptstyle -T$}} % 调整高度为 -0.85ex
\newcommand{\mone}{\raisebox{0.1ex}{$\scriptstyle -1$}} % 调整高度为 0.1ex
\newcommand{\lmone}{\raisebox{-0.85ex}{$\scriptstyle -1$}} % 调整高度为 -0.85ex
\newcommand{\lsup}[1]{\raisebox{-0.1ex}{$\scriptstyle #1$}}
\newcommand{\llsup}[1]{\raisebox{-0.85ex}{$\scriptstyle #1$}}
\newcommand{\mathminus}{\!\!-\!\!} % 数学环境连字符
\newcommand{\vvec}[1]{\symbf{#1}}
\newcommand{\ten}[1]{\mathbb{#1}}


\begin{document}
%\thispagestyle{empty}
\begin{center}
    \Huge{\textbf{我掌握的流体力学}}\\
    \vspace{1.5cm}
    \Huge\textbf{前言}
\end{center}

这份笔记整理始于2026年端午节，初衷是为推免申请做准备——希望在有限时间内，系统梳理流体力学
相关的核心概念，便于向老师系统展示我在本科前三年掌握的流体力学相关概念。在整理过程中，
复变函数、矢量恒等式等数学工具默认已在先修课程中掌握，故仅简要提及推导思路或者直接使用。
预计九月完成全部整理，涵盖流体力学基础（包括实验流体力学）、气体动力学、连续介质力学、理论力学，
以及热力学与统计物理等方向。目前正在撰写第一部分——流体力学基础。
\href{https://github.com/nobodyMa/nobodyMa_note_review_Fluiddynamics/blob/main/%E6%88%91%E6%8E%8C%E6%8F%A1%E7%9A%84%E6%B5%81%E4%BD%93%E5%8A%9B%E5%AD%A6.pdf}{最新版本的笔记}将持续更新在GitHub上，
欢迎大家提出建议。
\vspace{0.5cm}

\begin{flushright}
    作者：无名氏马\\
    \today
\end{flushright}

\newpage                      % 新的一页
\pagestyle{plain}             % 设置页眉和页脚的排版方式（plain：页眉是空的，页脚只包含一个居中的页码）
\setcounter{page}{1}          % 重新定义页码从第一页开始
\pagenumbering{Roman}         % 使用大写的罗马数字作为页码
\tableofcontents              % 生成目录

\newpage                      % 以下是正文
\pagestyle{plain}
\setcounter{page}{1}          % 使用阿拉伯数字作为页码
\pagenumbering{arabic}
% \setcounter{chapter}{-1}    % 设置 -1 可作为第零章绪论从第零章开始
 % 单独编译时，其实不用编译封面目录之类的，如需要不注释这句即可
\else
\fi
%  ↓↓↓↓↓↓↓↓↓↓↓↓↓↓↓↓↓↓↓↓↓↓↓↓↓↓↓↓ 正文部分
\chapter{流体力学基础}
\section{前置知识及基本概念}
\subsection{回忆}
数学工具：多元函数微积分（偏导数、全微分），矢量代数与矢量微积分（梯度、散度、旋度）
，奥--高定理，斯托克斯定理，二阶张量（对称张量、反对称张量、张量分解），赝矢量，
常微分方程，偏微分方程（拉普拉斯方程、扩散方程），量纲分析（白金汉\(\pi\) 定理）。

运动学基本概念：连续介质假设，流体微团（流体质点），欧拉描述，拉格朗日描述，
物质导数（随体导数），当地导数，对流项，流线，迹线（轨迹线），脉线（染色线）。
\subsection{应力张量} 
在连续介质力学中将流体视为连续填充空间的物质。此时，相邻流体部分之间的相互作用力通过假想截面上的面力分布来描述。

考虑流体内部一个以单位法向量 $\mathbf{n}$ 的假想面元 $\mathrm{d}S$。
令 $\mathrm{d}\mathbf{f}$ 为该面元两侧流体之间的相互作用力，可写成：
\begin{equation}
\mathrm{d}\mathbf{f} = \mathbf{t}(\mathbf{n}) \, \mathrm{d}S,
\end{equation}
其中 $\mathbf{t}(\mathbf{n})$ 称为应力矢量（或称牵引力矢量），它依赖于该点处面元的空间取向 $\mathbf{n}$。

\begin{principle}[柯西应力原理]
    应力矢量 $\mathbf{t}(\mathbf{n})$ 与法向量 $\mathbf{n}$ 之间的关系是线性的。存在一个二阶张量 $\boldsymbol{\sigma}$，使得：
    \begin{equation}
    \mathbf{t}(\mathbf{n}) = \boldsymbol{\sigma} \cdot \mathbf{n}\quad \text{或} \quad t_i(\mathbf{n}) = \sigma_{ij} n_j
     \end{equation}
     其中 $\boldsymbol{\sigma}$ 即为应力张量，其分量 $\sigma_{ij}$ 表示作用在法向为 $j$ 的面元上、沿 $i$ 方向的应力分量。
\end{principle}
\begin{proof}[量级分析]
取一个以 $O$ 为顶点、斜面法向为 $\mathbf{n}$ 的微小四面体，其三个坐标面分别垂直于
 $x_1, x_2, x_3$ 轴，面积分别为斜面面积\(\mathrm{d}S\)乘$\mathbf{n}$ 的分量 $n_1, n_2, n_3$。令四面体的特征尺度为 $h$，则体积 $\propto h^3$，面积 $\propto h^2$。
 当 $h \to 0$ 时，体积力（如重力、惯性力）的量级为 $O(h^3)$，远小于面力的量级 $O(h^2)$，
 因此体积力在极限下不影响平衡，四个面上的面力之和为零。
 \[
\mathbf{t}(\mathbf{n}) \, \mathrm{d}S + \mathbf{t}(-\mathbf{e}_1) \, \mathrm{d}S_1 + \mathbf{t}(-\mathbf{e}_2) \, \mathrm{d}S_2 + \mathbf{t}(-\mathbf{e}_3) \, \mathrm{d}S_3 = \mathbf{0}
\]
约去公共因子 \(\mathrm{d}S\)，得到
\[
\mathbf{t}(\mathbf{n}) = n_1 \mathbf{t}(\mathbf{e}_1) + n_2 \mathbf{t}(\mathbf{e}_2) + n_3 \mathbf{t}(\mathbf{e}_3)
\]
\end{proof}

\begin{proposition}
在\textcolor{red}{流体微元上的合力矩为零}条件下，应力张量是对称的：
\begin{equation}
\sigma_{ij} = \sigma_{ji}
\end{equation}
\begin{proof}[量级分析]
考虑一个边长分别为 \(\mathrm{d}x, \mathrm{d}y, \mathrm{d}z\) 的微小立方体流体微元。
对绕 \(z\) 轴的力矩，仅有切向分量 \(\sigma_{xy}\) 和 \(\sigma_{yx}\) 产生净贡献。
作用于与 \(x\) 轴垂直的两个面上的 \(\sigma_{xy}\) 构成一对反向力偶，其合力矩为；
作用于与 \(y\) 轴垂直的两个面上的 \(\sigma_{yx}\) 构成一对正向力偶，其合力矩为。
因此绕 \(z\) 轴的净力矩为
\[
M_z = (\sigma_{yx}-\sigma_{xy}) \, \mathrm{d}x \, \mathrm{d}y \, \mathrm{d}z,
\]
其量级为 \(O(h^3)\)。立方体的转动惯量 \(I \propto \rho h^5\)，角加速度 \(\dot{\omega}\) 的量级至多为 \(O(1)\)，
故惯性力矩 \(I\dot{\omega}\) 的量级为 \(O(h^5)\)。当微元尺寸趋于零时惯性力矩可忽略。为保持力矩平衡
\[
M_z = 0 \quad \Rightarrow \quad \sigma_{xy} = \sigma_{yx}
\]
\end{proof}
\end{proposition}
\begin{proposition}
    应力张量 $\boldsymbol{\sigma}$ 可以唯一地分解为一个球量部分（各项同性的静水压力）和一个偏量部分（流体黏性引起的各向异性应力）：
\begin{equation}\label{eq:stress_composition}
\boldsymbol{\sigma} = -p \mathbf{I} + \boldsymbol{\sigma}'\quad \text{或} \quad \sigma_{ij} = -p \delta_{ij} + \sigma'_{ij}
\end{equation}
其中：
\begin{itemize}
    \item $p = \frac{1}{3}\mathrm{tr}(\boldsymbol{\sigma}) = \frac{1}{3}\sigma_{kk}$ ，代表静水压力（压应力为正）；
    \item $\boldsymbol{\sigma}'$ 黏性应力张量，无迹 $\mathrm{tr}(\boldsymbol{\sigma}') = 0$，代表由流体黏性引起的各向异性应力，导致流体发生剪切变形；
    \item $\mathbf{I}$ 为单位二阶张量（克罗内克符号 $\delta_{ij}$）。
\end{itemize}
\end{proposition}
\begin{note}
    可以发现流体力学的应力张量是在流体静力学的基础上修正出的。
\end{note}

\subsection{流体的控制方程} 
\paragraph{连续性方程（质量守恒方程）}

考虑流体中任意一个\textcolor{red}{在参考系下固定的体积} $\mathcal{V}$，其边界为封闭曲面 $\mathcal{S}$，单位外法向量为 $\mathbf{n}$。
令 $\rho(\mathbf{r}, t)$ 为流体密度，$\mathbf{v}(\mathbf{r}, t)$ 为流速，则质量守恒的积分形式可写为：
\begin{equation}
\frac{d}{dt} \left( \iiint_{\mathcal{V}} \rho \, dV \right) 
=\iiint_{\mathcal{V}} \frac{\partial \rho}{\partial t} \, dV+0
= - \iint_{S} \rho \, \mathbf{v} \cdot \mathbf{n} \, dS
\end{equation}
对第二个等式应用奥--高定理：
\begin{equation}
\iiint_{\mathcal{V}} \left[ \frac{\partial \rho}{\partial t} + \nabla \cdot (\rho \mathbf{v}) \right] dV = 0
\end{equation}
由于体积 \( \mathcal{V} \) 的选择具有任意性，可得到连续方程的微分形式：
\begin{equation}\label{eq:continuity_euler}
\boxed{\frac{\partial \rho}{\partial t} + \nabla \cdot (\rho \mathbf{v}) = 0}
\end{equation}
此为欧拉观点下的连续性方程，因为选择的区域在描述流动的参考系下是固定的。

\begin{note}
类比，电磁学中的电荷守恒方程\(\frac{\partial \rho}{\partial t} + \nabla \cdot \mathbf{J} = 0\)，
其中\(\rho\)表示电荷密度，\(\mathbf{J}\)表示电流密度。
\end{note}

将方程\ref{eq:continuity_euler}的散度项展开，得到拉格朗日观点下的连续性方程：
\begin{equation}\label{eq:continuity_lagrange}
\frac{\partial \rho}{\partial t} + \rho \nabla \cdot \mathbf{v} + \mathbf{v} \cdot \nabla \rho
=\boxed{\frac{d\rho}{dt} + \rho \nabla \cdot \mathbf{v} = 0}
\end{equation}
该形式表明，对于跟随流体微团运动的观察者来说，
密度的变化率等于流体微团的体积膨胀率的负值 \( -\rho \nabla \cdot \mathbf{v} \)。

对于\textcolor{red}{不可压缩流体}，每个流体微团的密度沿其运动轨迹保持不变\(d\rho/dt = 0\)，
连续性方程简化为速度场的散度为零：
\begin{equation}\label{eq:continuity_incompressible}
\boxed{\nabla \cdot \mathbf{v} = 0}
\end{equation}

\paragraph{本构方程（牛顿流体）}
牛顿流体满足以下两条基本假设：
\begin{enumerate}
    \item 黏性应力与应变率张量成线性关系；
    \item 流体为各向同性介质。
\end{enumerate}

应变率张量定义为：
\begin{equation}
\mathbf{e} = \frac{1}{2} \left( \nabla \mathbf{v} + (\nabla \mathbf{v})^T \right), \quad e_{ij} = \frac{1}{2}\left( \frac{\partial v_i}{\partial x_j} + \frac{\partial v_j}{\partial x_i} \right).
\end{equation}

对于各向同性线性材料，最一般的四阶张量形式可简化为两个独立系数。
\begin{equation}\label{eq:constitutive_general}
\boldsymbol{\sigma}' = 2A\mathbf{e} + B (\nabla \cdot \mathbf{v})
\quad \text{或}\quad \sigma'_{ij} = 2Ae_{ij} + B \, \delta_{ij}e_{ll}
\end{equation}

由此可得：
\begin{equation}\label{eq:constitutive_newtonian}
\boldsymbol{\sigma}' = 2\eta \mathbf{e} + \left( \zeta - \frac{2}{3}\eta \right) (\nabla \cdot \mathbf{v}) \mathbf{I}
\quad \text{或}\quad 
\sigma'_{ij} = \eta \left(2e_{ij} - \frac{2}{3} \delta_{ij} e_{ll}\right) + \zeta \, \delta_{ij}e_{ll}
\end{equation}
其中 $\eta$ 为剪切黏度（动力黏度），$\zeta$ 为体积黏度（第二黏度）。第一项对应无体积变化时的纯剪切变形，第二项对应各向同性的胀缩变形。

\begin{explain}
类似于应力张量，将应变率张量分解为偏量部分（无迹）和球量部分（各向同性）：
\begin{equation}
\epsilon_{ij} = \underbrace{\left(\epsilon_{ij} - \frac{1}{3}\delta_{ij}\epsilon_{kk}\right)}_{\text{偏量}} + \underbrace{\frac{1}{3}\delta_{ij}\epsilon_{kk}}_{\text{球量}}.
\end{equation}

代入一般形式\ref{eq:constitutive_general}，依旧分解为偏量部分（无迹）和球量部分（各向同性）得：
\[
\begin{aligned}
\sigma'_{ij} &= A \epsilon_{ij} + B \delta_{ij}\epsilon_{kk} \\
&= A\left[\left(\epsilon_{ij} - \frac{1}{3}\delta_{ij}\epsilon_{kk}\right) + \frac{1}{3}\delta_{ij}\epsilon_{kk}\right] + B\delta_{ij}\epsilon_{kk} \\
&= A\left(\epsilon_{ij} - \frac{1}{3}\delta_{ij}\epsilon_{kk}\right) + \left(B + \frac{A}{3}\right)\delta_{ij}\epsilon_{kk}
\end{aligned}
\]

令剪切黏度 $\eta \equiv A$，体积黏度 $\zeta \equiv B + A/3$，则：
\[
\sigma'_{ij} = \eta \left(2e_{ij} - \frac{2}{3} \delta_{ij} e_{ll}\right) + \zeta \, \delta_{ij}e_{ll}
\]
\end{explain}

对于不可压缩流体，$\nabla \cdot \mathbf{v} = 0$，本构方程简化为：
\begin{equation}
\boldsymbol{\sigma}' = 2\eta \mathbf{e} = \eta \left( \nabla \mathbf{v} + (\nabla \mathbf{v})^T \right)
\quad \text{或} \quad
\sigma'_{ij} = \eta \left( \frac{\partial v_i}{\partial x_j} + \frac{\partial v_j}{\partial x_i} \right)
\end{equation}

\paragraph{流体运动方程（动量守恒方程）}
基于牛顿第二定律，一个\textcolor{red}{随流体运动的物质体积} $\mathcal{V}(t)$内流体的总动量变化率等于作用于其上的体积力和表面力之和：
\begin{equation}
\frac{\mathrm{d}}{\mathrm{d}t} \iiint_{\mathcal{V}(t)} \rho \mathbf{v} \,\mathrm{d}V 
= \iiint_{\mathcal{V}(t)} \rho \mathbf{f} \,\mathrm{d}V + \iint_{\partial \mathcal{V}(t)} \boldsymbol{\sigma} \cdot \mathbf{n} \,\mathrm{d}S
\end{equation}
其中 $\rho$ 为密度，$\mathbf{v}$ 为速度，$\mathbf{f}$ 为单位质量体积力，$\boldsymbol{\sigma}$ 为应力张量，$\mathbf{n}$ 为控制体表面的外法线单位向量。

左侧由于$\rho \mathrm{d}V$ 在随体运动（拉格朗日观点）中保持不变，右侧使用奥--高定理，
得到动量方程的积分形式：
\begin{equation}
\frac{\mathrm{d}}{\mathrm{d}t} \iiint_{\mathcal{V}(t)} \rho \mathbf{v} \,\mathrm{d}V 
= \boxed{\iiint_{\mathcal{V}(t)} \rho \frac{\mathrm{d}\mathbf{v}}{\mathrm{d}t} \,\mathrm{d}V
= \iiint_{\mathcal{V}(t)} \left(\rho \mathbf{f} + \nabla \cdot \boldsymbol{\sigma}\right) \,\mathrm{d}V
}\end{equation}

由于物质体积 $\mathcal{V}(t)$ 是任意的，得到局部的微分形式：
\begin{equation}
\rho \frac{\mathrm{d}\mathbf{v}}{\mathrm{d}t} = \rho \mathbf{f} + \nabla \cdot \boldsymbol{\sigma}.
\label{eq:motion_local}
\end{equation}

将应力张量分解式\ref{eq:stress_composition}入式\ref{eq:motion_local}，得到：
\begin{equation}
\rho \frac{\mathrm{d}\mathbf{v}}{\mathrm{d}t} = \rho \mathbf{f} - \nabla p + \nabla \cdot \boldsymbol{\sigma}'.
\label{eq:motion_local_pressure}
\end{equation}

上面是适用于任意连续介质的拉格朗日观点下运动方程，下面是欧拉观点下得到的运动方程。
将物质导数展开 $\mathrm{d}/\mathrm{d}t = \partial/\partial t + \mathbf{v}\cdot\nabla$，
得到欧拉观点下任意连续介质的运动方程：
\begin{equation}
\rho \left( \frac{\partial \mathbf{v}}{\partial t} + \mathbf{v}\cdot\nabla \mathbf{v} \right) = \rho \mathbf{f} - \nabla p + \nabla \cdot \boldsymbol{\sigma}'.
\label{eq:eulerian_general}
\end{equation}

将牛顿流体的本构方程\ref{eq:constitutive_newtonian}代入式\ref{eq:eulerian_general}，
得到纳维-斯托克斯方程：
\begin{equation}
\rho \left( \frac{\partial \mathbf{v}}{\partial t} + \mathbf{v}\cdot\nabla \mathbf{v} \right) = \rho \mathbf{f} - \nabla p + \eta \nabla^2 \mathbf{v} + \left( \zeta + \frac{\eta}{3} \right) \nabla (\nabla\cdot\mathbf{v}).
\label{eq:ns_compressible}
\end{equation}

若流体为不可压缩（$\nabla\cdot\mathbf{v}=0$），则体积黏度 $\zeta$ 项消失，
得不可压缩纳维-斯托克斯方程：
\begin{equation}
\rho \left( \frac{\partial \mathbf{v}}{\partial t} + \mathbf{v}\cdot\nabla \mathbf{v} \right) = \rho \mathbf{f} - \nabla p + \eta \nabla^2 \mathbf{v}.
\label{eq:ns_incompressible}
\end{equation}

当流体黏性可忽略时（$\eta = 0$，$\boldsymbol{\sigma}' = 0$），式\ref{eq:eulerian_general}退化为欧拉方程：
\begin{equation}
\rho \left( \frac{\partial \mathbf{v}}{\partial t} + \mathbf{v}\cdot\nabla \mathbf{v} \right) = \rho \mathbf{f} - \nabla p.
\label{eq:euler_incompressible}
\end{equation}

积分形式的运动方程详见\ref{sec:conservation}节守恒律。
下面将从欧拉观点下的微分形式运动方程\ref{eq:eulerian_general}出发，推导其散度形式。

利用连续性方程 $\partial\rho/\partial t + \nabla\cdot(\rho\mathbf{v}) = 0$，将左侧第一项展开：
\begin{equation}
\begin{aligned}
    \rho \frac{\partial \mathbf{v}}{\partial t} = \frac{\partial}{\partial t}(\rho \mathbf{v}) - \mathbf{v}\frac{\partial \rho}{\partial t}
= \frac{\partial}{\partial t}(\rho \mathbf{v}) + \mathbf{v}\,\nabla\cdot(\rho\mathbf{v}).
\end{aligned}
\end{equation}

将左侧第二项（对流项） $\rho(\mathbf{v}\cdot\nabla)\mathbf{v}$展开：
\begin{equation}
\rho(\mathbf{v}\cdot\nabla)\mathbf{v} 
= \rho v_j \frac{\partial v_i}{\partial x_j}
= \frac{\partial}{\partial x_j}(\rho v_i v_j) - v_i \frac{\partial}{\partial x_j}(\rho v_j)
= \nabla\cdot(\rho \mathbf{v}\otimes\mathbf{v}) - \mathbf{v}\,\nabla\cdot(\rho\mathbf{v})
\end{equation}
其中 $\mathbf{v}\otimes\mathbf{v}$ 表示并矢。

代入方程\ref{eq:eulerian_general}的左侧，得到：
\begin{equation}
\rho \frac{\partial \mathbf{v}}{\partial t} + \rho(\mathbf{v}\cdot\nabla)\mathbf{v}
= \frac{\partial}{\partial t}(\rho \mathbf{v}) + \nabla\cdot(\rho \mathbf{v}\otimes\mathbf{v})
= \rho \mathbf{f} - \nabla p + \nabla \cdot \boldsymbol{\sigma}'
\end{equation}

将所有含散度的项移到右侧，整理为守恒律形式：
\begin{equation}
\frac{\partial}{\partial t}(\rho \mathbf{v})
= -\nabla\cdot(\rho \mathbf{v}\otimes\mathbf{v} + p\mathbf{I} - \boldsymbol{\sigma}') + \rho \mathbf{f}
\quad \text{或} \quad
\frac{\partial}{\partial t}(\rho v_i) = -\frac{\partial}{\partial x_j}\left(\rho v_i v_j + p\delta_{ij} - \sigma'_{ij}\right) + \rho f_i
\end{equation}

\begin{definition}[动量通量张量 $\boldsymbol{\Pi}$]
    \begin{equation}
    \boldsymbol{\Pi} = \rho \mathbf{v}\otimes\mathbf{v} + p\mathbf{I} - \boldsymbol{\sigma}'
    \quad \text{或} \quad
    \Pi_{ij} = \rho v_i v_j + p\delta_{ij} - \sigma'_{ij}
    \end{equation}
其中各项的物理含义为：
\begin{itemize}
    \item $\rho v_i v_j$：由流体宏观运动携带的动量通量（对流项）；
    \item $p\delta_{ij}$：由压力传递的动量通量（各向同性挤压）；
    \item $-\sigma'_{ij}$：由黏性应力传递的动量通量（切向和法向黏性应力）。
\end{itemize}
\end{definition}

\paragraph{动能守恒方程}
不可压缩流体动能守恒方程、伯努利方程以及势流伯努利方程详见\ref{sec:conservation}节守恒律。

\paragraph{气体状态方程}
详见热力学与统计物理部分。
\subsection{无量纲数及其无量纲方程}
\paragraph{输运现象中的扩散与对流竞争}
\begin{table}[h]
\centering
\caption{输运现象中的扩散与对流竞争}
\begin{tabular}{|c|c|c|}
\hline
\textbf{无量纲数} & \textbf{定义} & \textbf{物理含义} \\
\hline
雷诺数 ($Re$) & $\dfrac{UL}{\nu}$ & 惯性力 / 黏性力；动量对流率 / 动量扩散率 \\
\hline
施密特数 ($Sc$) & $\dfrac{\nu}{D_m}$ & 动量扩散率 / 质量扩散率 \\
\hline
普朗特数 ($Pr$) & $\dfrac{\nu}{\kappa}$ & 动量扩散率 / 热量扩散率\\
\hline
刘易斯数 ($Le$) & $\dfrac{\kappa}{D_m}$ & 热量扩散率 / 质量扩散率\\
\hline
佩克莱数 ($Pe$) & $\dfrac{UL}{D_m} = Re \cdot Sc$ & 质量对流率 / 质量扩散率 \\
\hline
热佩克莱数 ($Pe_\theta$) & $\dfrac{UL}{\kappa} = Re \cdot Pr$ & 热量对流率 / 热量扩散率 \\
\hline
\end{tabular}
\end{table}























%  ↑↑↑↑↑↑↑↑↑↑↑↑↑↑↑↑↑↑↑↑↑↑↑↑↑↑↑↑ 正文部分
\ifx\allfiles\undefined
\end{document}
\fi